\documentclass[12pt]{article}

% --- Packages ---
\usepackage[margin=1in]{geometry}
\usepackage{amsmath,amssymb,amsthm}
\usepackage{mathtools}
\usepackage{graphicx}
\usepackage{booktabs}
\usepackage{natbib}
\usepackage{hyperref}
\usepackage{cleveref}
\usepackage{algorithm}
\usepackage{algpseudocode}

% --- Theorem environments ---
\newtheorem{theorem}{Theorem}
\newtheorem{lemma}[theorem]{Lemma}
\newtheorem{proposition}[theorem]{Proposition}
\newtheorem{corollary}[theorem]{Corollary}
\theoremstyle{definition}
\newtheorem{definition}[theorem]{Definition}
\theoremstyle{remark}
\newtheorem{remark}[theorem]{Remark}

% --- Notation shortcuts ---
\newcommand{\E}{\mathbb{E}}
\newcommand{\Var}{\mathrm{Var}}
\newcommand{\Rset}{\mathbb{R}}
\newcommand{\bone}{\mathbf{1}}
\newcommand{\blam}{\boldsymbol{\lambda}}
\DeclareMathOperator*{\argmax}{arg\,max}
\DeclareMathOperator*{\argmin}{arg\,min}

% --- Title ---
\title{Exponential Parallel System Estimation\\via the Censoring Equivalence}
\author{Alexander Towell\thanks{Department of Computer Science,
  Southern Illinois University Edwardsville.
  Email: \texttt{lex@metafunctor.com}.
  ORCID: 0000-0001-6443-9897.}}
\date{}

\begin{document}
\maketitle

\begin{abstract}
We show that observing the lifetime of a parallel system is equivalent to
observing one exact and $m-1$ left-censored component lifetimes with latent
identity. This censoring equivalence---formalized for independent exponential
components---enables closed-form maximum likelihood estimation via
inclusion-exclusion, an EM algorithm with analytical E- and M-steps, and an
elementary proof that the component rates are identifiable from system data
alone. The identifiability result stands in sharp contrast to the classical
non-identifiability of series (competing risks) systems, revealing a
structural duality: series systems mask \emph{which} component failed,
while parallel systems mask \emph{when} each component failed. Monte Carlo
simulations confirm that the MLE is consistent with nominal coverage and
reveal an information asymmetry governed by the probability that each
component is the last to fail: slow components are estimated precisely,
while fast components carry little system-level information.  A Fisher
information analysis quantifies this asymmetry, showing that the
per-component efficiency is approximately proportional to $P(T_j = T)$
and that the joint efficiency collapses exponentially with the number of
components.  The censoring equivalence extends to $k$-out-of-$n$ systems,
unifying the parallel ($k = 1$) and series ($k = n$) cases within a
single framework.
\end{abstract}

\noindent\textbf{Keywords:} parallel systems, $k$-out-of-$n$ systems,
exponential distribution, maximum likelihood, identifiability, censoring,
Fisher information, EM algorithm, reliability

% Section 1: Introduction
\section{Introduction}
\label{sec:introduction}

A parallel system of $m$ components functions as long as at least one
component survives; its lifetime is $T = \max(T_1, \ldots, T_m)$.  Such
systems arise in redundant engineering designs---backup power supplies,
server clusters, communication networks---where each component provides
an independent path to system functionality.  A basic engineering
question is: given observed system lifetimes, what can be learned about
the individual component lifetime distributions?

This question arises in practice whenever component-level testing is
impractical or impossible: satellites with redundant subsystems
observable only at system level, wind turbine gearboxes where
individual bearing instrumentation is costly, nuclear safety systems
tested by activating the entire redundant chain, and data center
servers whose component failures are recorded only at the machine
level.  In each case, system failure times are the primary data source,
and component-level inference requires exploiting the system structure.

For the dual topology---a series system with
$T = \min(T_1, \ldots, T_m)$---this question has generated a vast
literature under the heading of \emph{competing risks} and
\emph{masked cause of failure}.  The foundational result of
\citet{Tsiatis1975} establishes that component distributions are
\emph{not} identifiable from system data alone; additional
information about the cause of failure is needed.  Estimation under
partial cause information (masking) has been studied extensively
\citep{Miyakawa1984, UsherHodgson1988, LinUsherGuess1993,
GuessUsherHodgson1991, KunduBasu2000, CraiuReiser2006}.

The parallel case has received far less attention, despite sitting on
the opposite side of an identifiability divide.  \citet{BasuGhosh1980}
showed that the distribution of $\max(T_1, \ldots, T_m)$ uniquely
determines the component distributions under independence---a result
with no series counterpart.  Yet no paper has derived a closed-form
maximum likelihood estimator or EM algorithm for the component rates
of an exponential parallel system from system-level data.
\citet{SarhanElBassiouny2003} treat the exponential parallel system
under a masking framing, deriving MLE and Bayes estimators
numerically, but do not establish identifiability or exploit the
inclusion-exclusion structure.  \citet{CramerNavarro2015} show that
component failure times from a coherent system can be represented as
progressively Type-II censored order statistics---a result that
specializes to our censoring equivalence for the parallel case---but
do not derive closed-form exponential MLE from it.  Other approaches
are nonparametric and signature-based
\citep{BhattacharyaSamaniego2010, NgNavarroBalakrishnan2017,
YangNgBalakrishnan2019, Asadi2020}, Bayesian without closed-form
posteriors \citep{BheringPereiraPolpo2014}, or focused on
non-estimation aspects of system reliability.
Recent work has extended the problem in complementary directions:
\citet{QuNgMoon2024a,QuNgMoon2024b} develop nonparametric tests for
component distribution homogeneity from system data,
\citet{BiswasGangulyMitra2025} propose semiparametric models for
load-sharing parallel systems, and \citet{Jasinski2023,Jasinski2024}
derive closed-form MLE for discrete (geometric) lifetimes in
$k$-out-of-$n$ systems.  None of these addresses the continuous
exponential case with the censoring perspective developed here.

The key obstacle has been a misidentification of the incomplete-data
structure.  Several authors
\citep{LiuShiZhangBai2017, RodriguesPereiraBheringPolpo2018,
Sarhan2001} have transplanted the \emph{masking} framework
(conditions C1--C3) from series to parallel systems.  But masking
presupposes uncertainty about \emph{which} component caused system
failure---a question that is meaningful for series systems (where one
component fails and the rest survive) but structurally meaningless for
parallel systems (where every component has failed by the time the
system fails).

The natural incomplete-data problem for parallel systems is not masking
but \emph{censoring}.  When the system fails at time $T = t$, we know
that all components have failed and that exactly one failed at time~$t$
(the last survivor), but we do not know \emph{when} the other
components failed---only that they failed before~$t$.  This is
component-level left-censoring with latent identity, and it is the
correct framing for estimation.  The censoring perspective was
introduced by \citet{BheringPereiraPolpo2014} in a Bayesian Weibull
context, but its frequentist consequences have not been developed.

\subsection*{Contributions}

We formalize the censoring equivalence and exploit it for exponential
parallel systems with independent components.

\begin{enumerate}
  \item \textbf{Censoring equivalence} (\Cref{prop:censoring-equiv}).
    We prove that observing $T = \max(T_1, \ldots, T_m)$ is
    informationally equivalent to observing one exact and $m-1$
    left-censored component lifetimes with latent identity $J =
    \argmax_j T_j$.

  \item \textbf{Closed-form likelihood via inclusion-exclusion}
    (\Cref{lem:ie,lem:density-decomp,lem:integral-wj}).
    The system density decomposes as $f_T(t) = \sum_j w_j(t)$, where
    each component weight $w_j$ and its integrals have closed forms
    through the inclusion-exclusion expansion of the system CDF.
    This yields a tractable log-likelihood for exact, right-censored,
    and interval-censored system observations without numerical
    quadrature.

  \item \textbf{Elementary identifiability proof}
    (\Cref{thm:identifiability,cor:duality}).
    We provide a self-contained proof that the distribution of
    $T = \max T_j$ uniquely determines the component rates, by
    induction on the number of components via reversed hazard rate
    tail extraction.  A duality corollary formalizes the contrast:
    series systems determine only $\sum \lambda_j$; parallel systems
    determine all individual~$\lambda_j$.

  \item \textbf{EM algorithm with closed-form steps}
    (\Cref{prop:em-algorithm}).
    The censoring equivalence yields an EM algorithm whose E-step
    uses the posterior $P(J = j \mid T = t)$ and the truncated
    exponential mean, and whose M-step is a simple reciprocal.
    Both steps are in closed form---no simulation or numerical
    integration is needed.

  \item \textbf{Monte Carlo validation and information efficiency}
    (\Cref{sec:simulations}).
    Simulations confirm consistency and nominal coverage, and reveal
    an information asymmetry governed by $P(T_j = T)$: slow components
    are estimated precisely, while fast components carry little
    system-level information.  Fisher information analysis
    quantifies this: the per-component efficiency is approximately
    proportional to $P(T_j = T)$, and the joint efficiency
    collapses exponentially with~$m$.

  \item \textbf{$k$-out-of-$n$ generalization}
    (\Cref{prop:kofn-censoring}).
    The censoring equivalence extends to $k$-out-of-$n$ systems:
    observing $T = T_{(n-k+1)}$ yields one exact, $n-k$ left-censored,
    and $k-1$ right-censored component observations with latent identity.
    This unifies the parallel ($k=1$) and series ($k=n$) cases and
    reveals a monotone information structure as~$k$ varies.
\end{enumerate}

\subsection*{Outline}

\Cref{sec:model} introduces the model and the censoring equivalence.
\Cref{sec:likelihood} develops the inclusion-exclusion machinery and
establishes identifiability.
\Cref{sec:estimation} presents the direct MLE and EM algorithm.
\Cref{sec:simulations} reports simulation results, including a Fisher
information comparison.
\Cref{sec:discussion} extends the censoring equivalence to
$k$-out-of-$n$ systems and discusses limitations and future directions.

% Section 2: Model and Censoring Equivalence
\section{Model and Censoring Equivalence}
\label{sec:model}

\subsection{Model definition}
\label{sec:model-def}

Consider a parallel system of $m$ independent components with lifetimes
$T_1, \ldots, T_m$, where each component $j$ has an exponential distribution
with rate $\lambda_j > 0$:
\begin{equation}\label{eq:component-dist}
  T_j \sim \mathrm{Exp}(\lambda_j), \qquad
  F_j(t) = 1 - e^{-\lambda_j t}, \qquad
  f_j(t) = \lambda_j e^{-\lambda_j t}, \qquad t > 0.
\end{equation}
The system functions as long as at least one component survives, so the
system lifetime is
\begin{equation}\label{eq:system-lifetime}
  T = \max(T_1, \ldots, T_m).
\end{equation}
The system CDF and density are
\begin{align}
  F_T(t) &= \prod_{j=1}^m F_j(t) = \prod_{j=1}^m \bigl(1 - e^{-\lambda_j t}\bigr),
  \label{eq:system-cdf} \\
  f_T(t) &= \frac{d}{dt} F_T(t) = \sum_{j=1}^m f_j(t) \prod_{i \neq j} F_i(t),
  \label{eq:system-density}
\end{align}
the latter following from the product rule.  We write
$\blam = (\lambda_1, \ldots, \lambda_m) \in \Rset_{>0}^m$ for the parameter
vector.

Under \emph{Scheme~0} (black-box observation), we observe $n$ independent
copies $T^{(1)}, \ldots, T^{(n)}$ of the system lifetime~$T$.  No
component-level information is available: the observer sees only when the
system fails, not which components failed or in what order.  This is the
minimal observation scheme and the focus of the present paper.  More
informative schemes---periodic inspection, partial monitoring, continuous
component monitoring---are discussed briefly in \Cref{sec:discussion}.

\subsection{What system data reveals about components}
\label{sec:what-data-reveals}

When a system failure is observed at time $T = t$, the following facts
are logically entailed by the parallel structure~\eqref{eq:system-lifetime}:

\begin{enumerate}
  \item[(i)] \emph{All components have failed}: $T_j \leq t$ for every
    $j = 1, \ldots, m$.  Unlike the series system, where surviving
    components remain operational at system failure, the parallel system
    fails only after the last component expires.

  \item[(ii)] \emph{Exactly one component failed at time~$t$}: since the
    $T_j$ are independent and continuous, the probability of ties is zero,
    so almost surely there exists a unique index
    $J = \argmax_{j} T_j$ with $T_J = t$.

  \item[(iii)] \emph{The identity~$J$ is unobserved}: the system-level
    datum~$T = t$ does not reveal which component was the last to fail.
\end{enumerate}

\noindent
Each observation of $T$ carries the value of one component
lifetime exactly---but we do not know which one---and constrains all other
component lifetimes to lie below~$t$.  This structure is formalized in the
next subsection.

\subsection{The censoring equivalence}
\label{sec:censoring-equiv}

The following proposition recasts the incomplete-data structure of a parallel
system observation in the language of censoring.

\begin{proposition}[Censoring equivalence]\label{prop:censoring-equiv}
  The likelihood function of $\blam$ based on the system observation
  $T = \max(T_1, \ldots, T_m)$ is identical to the likelihood arising
  from the following component-level censoring scheme:
  \begin{enumerate}
    \item[\textup{(a)}] One component---the last to fail,
      $J = \argmax_j T_j$---is observed exactly: $T_J = T$.
    \item[\textup{(b)}] All other components are left-censored at~$T$: for
      $j \neq J$, we know only that $T_j < T$.
    \item[\textup{(c)}] The identity~$J$ is latent (unobserved).
  \end{enumerate}
\end{proposition}

\begin{proof}
Since $T = \max_j T_j$ and the component lifetimes are independent and
continuous, exactly one component satisfies $T_j = T$ almost surely; call
it~$J$.  For every other component, $T_j < T_J = T$.  The system-level
observation $T$ therefore carries precisely the following information:
\begin{itemize}
  \item the event $\{T_J = T\}$: one exact component observation;
  \item the events $\{T_j < T\}$ for all $j \neq J$: $m-1$ left-censored
    component observations at the common censoring time~$T$;
  \item the identity~$J$ is not observed.
\end{itemize}
No additional information is contained in the system observation $T = t$
beyond these three items, since $T = \max_j T_j$ is a deterministic function
of $(T_1, \ldots, T_m)$.  The complete data $(T_1, \ldots, T_m, J)$ augments
the observed data~$T$ with the latent identity and the left-censored lifetimes,
yielding the incomplete-data structure described in~(a)--(c).
\end{proof}

\Cref{prop:censoring-equiv} provides the conceptual foundation for the
estimation approach developed in \Cref{sec:likelihood,sec:estimation}.
The latent variable~$J$ plays a role analogous to the masked cause of
failure in series systems, but the censoring mechanism is fundamentally
different: left-censoring rather than right-censoring.
The result can also be obtained as a special case of the progressive
Type-II censoring representation for coherent systems
\citep{CramerNavarro2015}, which establishes that component failure times
from a system with known structure correspond to progressively censored
order statistics.  For a parallel system, this specializes to exactly
the one-exact plus $(m-1)$-left-censored structure of~(a)--(c).

\subsection{Contrast with series systems}
\label{sec:series-contrast}

\begin{remark}[Series--parallel duality in censoring]\label{rmk:series-parallel}
For a series system $T = \min(T_1, \ldots, T_m)$, the censoring equivalence
takes the dual form:
\begin{enumerate}
  \item[\textup{(a$'$)}] One component---the first to fail,
    $J = \argmin_j T_j$---is observed exactly: $T_J = T$.
  \item[\textup{(b$'$)}] All other components are \emph{right}-censored at~$T$:
    for $j \neq J$, we know only that $T_j > T$.
  \item[\textup{(c$'$)}] The identity~$J$ may be masked (observed, partially
    observed, or unobserved).
\end{enumerate}
The structural duality is: left-censoring in parallel systems corresponds to
right-censoring in series systems.  This mirrors the physical difference: in a
series system, surviving components have $T_j > T$ (they outlast the system); in
a parallel system, already-failed components have $T_j < T$ (they predecease the
system).

The notion of \emph{masking}---the conditions C1--C3 of
\citet{Miyakawa1984} and \citet{LinUsherGuess1993}---is physically motivated
for series systems, where the cause of system failure (which component failed
first) may be uncertain.  For parallel systems, masking is structurally
unmotivated: by the time the system fails, every component has failed, and
there is no single ``cause'' to mask.  Some authors
\citep{LiuShiZhangBai2017, RodriguesPereiraBheringPolpo2018} have
applied the C1--C3 framework to parallel systems, but the incomplete-data
problem is more naturally cast as component-level censoring with latent
identity, as in \Cref{prop:censoring-equiv}.  The censoring framing was
introduced by \citet{BheringPereiraPolpo2014} in a Bayesian Weibull setting;
we develop it here for frequentist exponential inference.
\end{remark}

\subsection{System-level censoring}
\label{sec:system-censoring}

In practice, the system lifetime~$T$ may itself be subject to censoring.
\Cref{prop:censoring-equiv} describes the component-level incomplete-data
structure \emph{given an exact observation of~$T$}.  When $T$ is censored,
the component-level implications compound as follows.

\begin{center}
\begin{tabular}{lll}
  \toprule
  \textbf{System observation} & \textbf{Component implication} & \textbf{Likelihood term} \\
  \midrule
  Exact: $T = t$ & $\exists!\, j\!: T_j = t$; $T_i < t\;\forall i \neq j$ &
    $f_T(t; \blam)$ \\[4pt]
  Right-censored: $T > \tau$ & $\exists\, j\!: T_j > \tau$; others unknown &
    $1 - F_T(\tau; \blam)$ \\[4pt]
  Interval: $a < T \leq b$ & All $T_j \leq b$; $\exists\, j\!: T_j > a$ &
    $F_T(b; \blam) - F_T(a; \blam)$ \\
  \bottomrule
\end{tabular}
\end{center}

\noindent
The right-censored case (system still functioning at monitoring time~$\tau$)
is the most common in practice.  Here, at least one component survives past
$\tau$, but we do not know which, nor do we know the status of the remaining
components: some may have failed silently before~$\tau$.  The system-level
survival probability $1 - F_T(\tau; \blam) = 1 - \prod_j (1 - e^{-\lambda_j \tau})$
has a closed form via the inclusion-exclusion expansion.

The interval-censored case arises under periodic inspection
(Scheme~1): the system fails between inspection times $a$ and $b$, so
$a < T \leq b$.  The likelihood contribution
$F_T(b; \blam) - F_T(a; \blam)$ is also available in closed form.

All three observation types---exact, right-censored, and
interval-censored---are accommodated by the likelihood framework of
\Cref{sec:likelihood}, where the relevant closed-form
expressions are developed in detail.

% Section 3: Likelihood and Identifiability
\section{Likelihood and Identifiability}
\label{sec:likelihood}

The system CDF~\eqref{eq:system-cdf} is a product of terms
$(1 - e^{-\lambda_j t})$, each of which is itself a difference of
exponentials.  Expanding this product via inclusion-exclusion yields the
analytical machinery that underlies all subsequent results: closed-form
densities, tractable likelihoods, and an elementary identifiability proof.

\subsection{Inclusion-exclusion expansion}
\label{sec:ie-expansion}

\begin{lemma}[Inclusion-exclusion expansion]\label{lem:ie}
  For rates $\lambda_1, \ldots, \lambda_m > 0$ and $t > 0$,
  \begin{equation}\label{eq:ie-expansion}
    \prod_{j=1}^m \bigl(1 - e^{-\lambda_j t}\bigr)
    = \sum_{S \subseteq \{1,\ldots,m\}} (-1)^{|S|}\, e^{-\lambda(S)\, t},
  \end{equation}
  where $\lambda(S) = \sum_{j \in S} \lambda_j$, with the convention
  $\lambda(\varnothing) = 0$.
\end{lemma}

\begin{proof}
By induction on~$m$; see \Cref{sec:appendix-proofs}.
\end{proof}

\begin{remark}\label{rmk:ie-terms}
The expansion~\eqref{eq:ie-expansion} has $2^m$ terms.  Each term
$e^{-\lambda(S) t}$ is an exponential with rate equal to a subset sum of
the~$\lambda_j$.  When the rates are distinct and no two subset sums
coincide---a generic condition---all $2^m$ terms have distinct rates.
This structural richness is ultimately what makes parallel systems
identifiable (\Cref{thm:identifiability}).
\end{remark}

\subsection{Density decomposition}
\label{sec:density-decomp}

\begin{lemma}[Density decomposition]\label{lem:density-decomp}
  The system density decomposes as
  \begin{equation}\label{eq:density-decomp}
    f_T(t; \blam) = \sum_{j=1}^m w_j(t; \blam),
  \end{equation}
  where the \emph{component weight function}
  \begin{equation}\label{eq:wj-def}
    w_j(t; \blam) = \lambda_j e^{-\lambda_j t} \prod_{i \neq j}
    \bigl(1 - e^{-\lambda_i t}\bigr)
  \end{equation}
  satisfies $w_j(t; \blam) \geq 0$ for all $t > 0$.  Moreover,
  \begin{equation}\label{eq:posterior-J}
    \frac{w_j(t; \blam)}{f_T(t; \blam)} = P(J = j \mid T = t; \blam),
  \end{equation}
  the conditional probability that component~$j$ was the last to fail,
  given system failure at time~$t$.
\end{lemma}

\begin{proof}
By direct computation of $P(T \in dt,\; J = j)$ via independence;
see \Cref{sec:appendix-proofs}.
\end{proof}

\begin{remark}\label{rmk:series-contrast-density}
The decomposition~\eqref{eq:density-decomp} is the parallel analogue
of the series decomposition $f_T(t) = \sum_j h_j(t) S_T(t)$, where
$h_j$ is the cause-specific hazard rate and $S_T$ is the system survival
function.  In the series case, hazard rate additivity
$h_T(t) = \sum_j h_j(t)$ provides a clean factorization.  The parallel
case has no such additivity; instead, the density is a sum of products,
with each term coupling one component's density to the CDFs of all
others.
\end{remark}

\subsection{Closed-form integrals}
\label{sec:closed-form-integral}

The weight functions $w_j$ can be integrated in closed form---a property
that is essential for handling censored and interval-censored system
observations without numerical quadrature.

\begin{lemma}[Closed-form integral of $w_j$]\label{lem:integral-wj}
  For $0 \leq a < b \leq \infty$ and all $\lambda_j > 0$,
  \begin{equation}\label{eq:integral-wj}
    \int_a^b w_j(t; \blam)\, dt
    = \lambda_j \sum_{S \subseteq \{1,\ldots,m\} \setminus \{j\}}
      (-1)^{|S|} \cdot
      \frac{e^{-r_S a} - e^{-r_S b}}{r_S},
  \end{equation}
  where $r_S = \lambda_j + \lambda(S) > 0$.
  In particular, when $b = \infty$,
  \begin{equation}\label{eq:integral-wj-tail}
    \int_a^\infty w_j(t; \blam)\, dt
    = \lambda_j \sum_{S \subseteq \{1,\ldots,m\} \setminus \{j\}}
      (-1)^{|S|} \cdot \frac{e^{-r_S a}}{r_S},
  \end{equation}
  and when $a = 0$ and $b = \infty$,
  \begin{equation}\label{eq:prob-last}
    P(J = j; \blam) = \int_0^\infty w_j(t; \blam)\, dt
    = \lambda_j \sum_{S \subseteq \{1,\ldots,m\} \setminus \{j\}}
      \frac{(-1)^{|S|}}{r_S},
  \end{equation}
  the unconditional probability that component~$j$ is the last to fail.
\end{lemma}

\begin{proof}
Apply \Cref{lem:ie} to the $m-1$ rates $\{\lambda_i : i \neq j\}$
appearing in the product within $w_j$:
\[
  w_j(t; \blam)
  = \lambda_j e^{-\lambda_j t} \prod_{i \neq j}
    \bigl(1 - e^{-\lambda_i t}\bigr)
  = \lambda_j \sum_{S \subseteq \{1,\ldots,m\} \setminus \{j\}}
    (-1)^{|S|}\, e^{-({\lambda_j + \lambda(S)})\, t}
  = \lambda_j \sum_S (-1)^{|S|}\, e^{-r_S t}.
\]
Since $r_S = \lambda_j + \lambda(S) \geq \lambda_j > 0$, each exponential
term is integrable on $[a, b]$:
\[
  \int_a^b e^{-r_S t}\, dt = \frac{e^{-r_S a} - e^{-r_S b}}{r_S}.
\]
Linearity of integration gives~\eqref{eq:integral-wj}.  The special
cases~\eqref{eq:integral-wj-tail} and~\eqref{eq:prob-last} follow by
substituting $b = \infty$ (so $e^{-r_S b} = 0$) and then $a = 0$
(so $e^{-r_S a} = 1$).
\end{proof}

\begin{remark}\label{rmk:integral-tractability}
\Cref{lem:integral-wj} is the key computational enabler for censored
likelihoods.  For right-censored system observations ($T > \tau$), the
survival probability $1 - F_T(\tau) = \int_\tau^\infty f_T(t)\, dt =
\sum_j \int_\tau^\infty w_j(t)\, dt$ reduces to a finite sum of
exponentials via~\eqref{eq:integral-wj-tail}.  For interval-censored
observations ($a < T \leq b$), the probability
$F_T(b) - F_T(a) = \sum_j \int_a^b w_j(t)\, dt$ likewise has a closed
form via~\eqref{eq:integral-wj}.  No numerical quadrature is needed.
\end{remark}

\subsection{Log-likelihood}
\label{sec:log-likelihood}

Suppose we observe $n$ system lifetimes, some exact and some censored.
Let $\mathcal{E}$, $\mathcal{R}$, and $\mathcal{I}$ denote the index
sets for exact, right-censored, and interval-censored observations,
respectively.  The observed-data log-likelihood is
\begin{equation}\label{eq:loglik}
  \ell(\blam) = \sum_{i \in \mathcal{E}} \ln f_T(t_i; \blam)
  + \sum_{i \in \mathcal{R}} \ln \bigl[1 - F_T(\tau_i; \blam)\bigr]
  + \sum_{i \in \mathcal{I}} \ln \bigl[F_T(b_i; \blam) - F_T(a_i; \blam)\bigr],
\end{equation}
where each term has a closed-form expression via the inclusion-exclusion
expansion:
\begin{align}
  f_T(t; \blam) &= \sum_{j=1}^m w_j(t; \blam)
  = \sum_{j=1}^m \lambda_j \sum_{S \subseteq \{1,\ldots,m\} \setminus \{j\}}
    (-1)^{|S|}\, e^{-r_S t},
  \label{eq:fT-ie} \\
  F_T(t; \blam) &= \sum_{S \subseteq \{1,\ldots,m\}}
    (-1)^{|S|}\, e^{-\lambda(S)\, t},
  \label{eq:FT-ie} \\
  F_T(b) - F_T(a) &= \sum_{S \subseteq \{1,\ldots,m\}}
    (-1)^{|S|}\, \bigl(e^{-\lambda(S)\, a} - e^{-\lambda(S)\, b}\bigr).
  \label{eq:interval-ie}
\end{align}

\noindent
Each evaluation of $\ell(\blam)$ requires $O(n \cdot m \cdot 2^{m-1})$
arithmetic operations.  The exponential scaling in~$m$ limits direct
application to systems with a moderate number of components
(say $m \leq 15$--$20$), but this covers many practical redundant
systems.

When all $n$ observations are exact ($\mathcal{R} = \mathcal{I} = \varnothing$),
the log-likelihood reduces to
\begin{equation}\label{eq:loglik-exact}
  \ell(\blam) = \sum_{i=1}^n \ln\!\Biggl[\,
    \sum_{j=1}^m w_j(t_i; \blam) \Biggr],
\end{equation}
the Scheme~0 (black-box) log-likelihood.

\subsection{Identifiability}
\label{sec:identifiability}

The central theoretical result of this paper is that the distribution of
the parallel system lifetime uniquely determines all component rates.

\begin{theorem}[Identifiability of exponential parallel systems]
\label{thm:identifiability}
  Let $T_1, \ldots, T_m$ be independent with $T_j \sim \mathrm{Exp}(\lambda_j)$,
  $\lambda_j > 0$, and let $T = \max(T_1, \ldots, T_m)$.  The distribution
  of~$T$ uniquely determines the multiset $\{\lambda_1, \ldots, \lambda_m\}$.

  That is, if $\blam = (\lambda_1, \ldots, \lambda_m)$ and
  $\boldsymbol{\mu} = (\mu_1, \ldots, \mu_m)$ both generate the same
  distribution for $T = \max_j T_j$, then
  $\{\lambda_1, \ldots, \lambda_m\} = \{\mu_1, \ldots, \mu_m\}$ as
  multisets.
\end{theorem}

\begin{proof}
The proof proceeds by induction on~$m$, extracting component rates one
at a time from the tail behavior of the reversed hazard rate.

\emph{Base case} ($m = 1$).
The CDF $F_T(t) = 1 - e^{-\lambda_1 t}$ determines $\lambda_1$
uniquely.

\emph{Inductive step.}  Assume the result holds for all parallel systems
with fewer than~$m$ components.  We show that $F_T$ determines the
multiset of rates.

\textbf{Step~1: Extract the smallest rate from the tail.}
The reversed hazard rate of~$T$
\citep[the natural hazard analogue for parallel systems;
see][]{BlockSavitsSingh1998} is
\[
  r_T(t) = \frac{f_T(t)}{F_T(t)}
  = \frac{d}{dt} \ln F_T(t)
  = \sum_{j=1}^m \frac{\lambda_j e^{-\lambda_j t}}{1 - e^{-\lambda_j t}}.
\]
For large~$t$, each summand admits the geometric series expansion
\[
  \frac{\lambda_j e^{-\lambda_j t}}{1 - e^{-\lambda_j t}}
  = \lambda_j e^{-\lambda_j t} \sum_{k=0}^\infty e^{-k\lambda_j t}
  = \lambda_j e^{-\lambda_j t} + O\bigl(e^{-2\lambda_j t}\bigr).
\]
Let $\lambda_* = \min_j \lambda_j$ denote the smallest rate, with
multiplicity~$\kappa$ (the number of components having rate~$\lambda_*$).
As $t \to \infty$,
\begin{equation}\label{eq:rhr-asymp}
  r_T(t) = \kappa\, \lambda_*\, e^{-\lambda_* t}
  + O\bigl(e^{-(\lambda_* + \delta) t}\bigr),
\end{equation}
where $\delta = \min\{\lambda_j - \lambda_* : \lambda_j > \lambda_*\} > 0$
if $\kappa < m$, and the remainder is exponentially smaller than the
leading term.  Since $r_T(t)$ is determined by $F_T$, we extract
$\lambda_*$ as
\[
  \lambda_* = \lim_{t \to \infty} \frac{-\ln r_T(t)}{t},
\]
and the multiplicity as
\[
  \kappa = \frac{1}{\lambda_*} \lim_{t \to \infty} r_T(t)\, e^{\lambda_* t}.
\]
Both limits exist and are uniquely determined by~$F_T$.

\textbf{Step~2: Factor out the identified components.}
Define
\[
  G(t) = \frac{F_T(t)}{\bigl(1 - e^{-\lambda_* t}\bigr)^\kappa}.
\]
Since $F_T(t) = (1 - e^{-\lambda_* t})^\kappa \prod_{j:\, \lambda_j > \lambda_*}
(1 - e^{-\lambda_j t})$, we have
\[
  G(t) = \prod_{j:\, \lambda_j > \lambda_*} \bigl(1 - e^{-\lambda_j t}\bigr),
\]
which is the CDF of $\max\{T_j : \lambda_j > \lambda_*\}$---a parallel system
with $m - \kappa < m$ components.  Since $F_T$ and $\lambda_*$ are known, $G$
is known.

\textbf{Step~3: Apply the induction hypothesis.}
By the induction hypothesis, the multiset $\{\lambda_j : \lambda_j > \lambda_*\}$
is uniquely determined by~$G$.  Combined with the $\kappa$ copies of~$\lambda_*$
identified in Step~1, the full multiset $\{\lambda_1, \ldots, \lambda_m\}$ is
determined.
\end{proof}

\begin{remark}\label{rmk:identifiability-role-of-independence}
The proof uses independence in two places: to factor $F_T$ as a product and
to decompose the reversed hazard rate as a sum.  Under dependence (e.g.,
copula models), identifiability may fail; see \citet{BasuGhosh1978} for
conditions and counterexamples.
\end{remark}

\subsection{Identifiability duality}
\label{sec:id-duality}

\begin{corollary}[Identifiability duality]\label{cor:duality}
  For independent exponential components with rates
  $\lambda_1, \ldots, \lambda_m > 0$:
  \begin{enumerate}
    \item[\textup{(a)}] \textbf{Series} ($T = \min_j T_j$).  The distribution
      of~$T$ determines only the total rate
      $\Lambda = \sum_{j=1}^m \lambda_j$, not the individual rates.
    \item[\textup{(b)}] \textbf{Parallel} ($T = \max_j T_j$).  The distribution
      of~$T$ determines the individual rates
      $\{\lambda_1, \ldots, \lambda_m\}$ without any additional information.
  \end{enumerate}
\end{corollary}

\begin{proof}
(a) The series survival function is
$\bar{F}_T(t) = \prod_j e^{-\lambda_j t} = e^{-\Lambda t}$, which depends
on~$\blam$ only through~$\Lambda$.  Any partition of~$\Lambda$ into~$m$
positive parts yields an observationally equivalent model.  This is
the classical non-identifiability result of \citet{Tsiatis1975}.

(b) \Cref{thm:identifiability}.
\end{proof}

\begin{remark}\label{rmk:duality-information}
The duality extends to the information-theoretic level.  The Fisher
information from system data alone has rank~$1$ for series systems (all
information concentrates in the direction $(1, \ldots, 1)$ in parameter
space) and full rank for parallel systems (\Cref{thm:fisher}).  The
incomplete-data problems are structural duals: series systems need
cause-of-failure data to separate the~$\lambda_j$; parallel systems are
already identified, but left-censoring of fast components degrades
estimation efficiency.
\end{remark}

\subsection{Fisher information}
\label{sec:fisher}

\begin{theorem}[Positive definiteness of Fisher information]
\label{thm:fisher}
  Under Scheme~0, the Fisher information matrix for a single observation
  of $T = \max(T_1, \ldots, T_m)$ with $T_j \sim \mathrm{Exp}(\lambda_j)$,
  \begin{equation}\label{eq:fisher-def}
    I(\blam)_{jk} = \E\!\left[
      \frac{\partial \ln f_T(T; \blam)}{\partial \lambda_j}
      \cdot \frac{\partial \ln f_T(T; \blam)}{\partial \lambda_k}
    \right],
  \end{equation}
  is positive definite for all $\blam \in \Rset_{>0}^m$ with distinct
  rates.
\end{theorem}

\begin{proof}
The Fisher information matrix is the variance-covariance matrix of the
score vector $\nabla_{\blam} \ln f_T(T; \blam)$, hence positive
semidefinite.  We must show it is nonsingular.

Suppose $I(\blam) v = 0$ for some $v \in \Rset^m$.  Then
\[
  \Var\!\bigl(v^\top \nabla_{\blam} \ln f_T(T; \blam)\bigr)
  = v^\top I(\blam)\, v = 0,
\]
so $v^\top \nabla_{\blam} \ln f_T(T; \blam) = 0$ almost surely.  That
is,
\[
  \sum_{j=1}^m v_j \frac{\partial}{\partial \lambda_j}
  \ln f_T(t; \blam) = 0
  \qquad \text{for almost every } t > 0,
\]
where the almost-sure statement under~$P$ transfers to
Lebesgue-almost-every~$t$ because $f_T(t; \blam) > 0$ for all $t > 0$.
This implies that the model is locally
non-identifiable in the direction~$v$.

By \Cref{thm:identifiability}, the model
$\{f_T(\cdot; \blam) : \blam \in \Rset_{>0}^m\}$ is globally
identifiable up to permutation of the rates.  For distinct rates, the
permutation equivalence class $\{\blam_\sigma : \sigma \in S_m\}$ is a
finite (hence discrete) set, through which no smooth curve can pass.
Therefore there are no
nontrivial directions of local non-identifiability, and $v = 0$.
\end{proof}

\begin{remark}\label{rmk:fisher-explicit}
While the positive definiteness of $I(\blam)$ is guaranteed by
\Cref{thm:fisher}, closed-form expressions for its entries require
integration of products of score components against $f_T$---integrals
that do not simplify to elementary form.  In practice, $I(\blam)$ is
computed numerically, either by Monte Carlo integration or by the
observed information matrix $-\nabla^2 \ell(\hat{\blam})$ evaluated
at the MLE.

The partial derivatives needed for the score are available analytically
via the inclusion-exclusion expansion.  From~\eqref{eq:wj-def},
differentiating $w_j$ with respect to $\lambda_j$ and $w_k$ ($k \neq j$)
with respect to $\lambda_j$ gives:
\begin{align}
  \frac{\partial w_j(t)}{\partial \lambda_j}
  &= \sum_{S \subseteq \{1,\ldots,m\} \setminus \{j\}}
    (-1)^{|S|}\, \bigl(1 - \lambda_j t\bigr)\, e^{-r_S t},
  \label{eq:dwj-dlamj} \\[6pt]
  \frac{\partial w_k(t)}{\partial \lambda_j}
  &= -\lambda_k \sum_{\substack{S \subseteq \{1,\ldots,m\} \setminus \{k\} \\
    j \in S}} (-1)^{|S|}\, t\, e^{-({\lambda_k + \lambda(S)})\, t},
  \qquad k \neq j.
  \label{eq:dwk-dlamj}
\end{align}
These are finite sums of exponentials, so the score and Hessian of
$\ell(\blam)$ are analytically tractable even though the expected Fisher
information is not.
\end{remark}

% Section 4: Maximum Likelihood Estimation
\section{Maximum Likelihood Estimation}
\label{sec:estimation}

Given $n$ independent observations $t_1,\ldots,t_n$ from the parallel
system $T = \max(T_1,\ldots,T_m)$ with $T_j \sim \mathrm{Exp}(\lambda_j)$,
the observed-data log-likelihood under Scheme~0 is
\eqref{eq:loglik-exact}, restated here for reference:
\[
  \ell(\blam)
  = \sum_{i=1}^{n} \ln f_T(t_i;\blam)
  = \sum_{i=1}^{n} \ln \Biggl[\,\sum_{j=1}^{m} w_j(t_i;\blam)\Biggr].
\]
\Cref{thm:identifiability} guarantees that the maximizer is unique
(up to permutation of equal rates), and \Cref{thm:fisher} ensures
that the Fisher information is positive definite at distinct rates, so the
MLE is well-defined.

We present two approaches to computing the MLE: direct numerical
maximization (\S\ref{sub:direct-mle}) and an EM algorithm with
closed-form steps (\S\ref{sub:em-algorithm}).  Both converge to the
same estimate; the EM is particularly attractive because it avoids
gradient computation and is guaranteed to increase the likelihood at
every iteration.

% ===================================================================
\subsection{Direct maximization}
\label{sub:direct-mle}
% ===================================================================

The log-likelihood is a smooth function of
$\blam \in \Rset_{>0}^m$, and its gradient and Hessian are available
analytically through the inclusion-exclusion expansion of
\Cref{lem:ie}.  The score vector has components
\begin{equation}
  \label{eq:score}
  \frac{\partial \ell}{\partial \lambda_j}
  = \sum_{i=1}^{n}
    \frac{1}{f_T(t_i)}
    \Biggl[
      \frac{\partial w_j(t_i)}{\partial \lambda_j}
      + \sum_{k \neq j} \frac{\partial w_k(t_i)}{\partial \lambda_j}
    \Biggr],
\end{equation}
where the partial derivatives~\eqref{eq:dwj-dlamj}--\eqref{eq:dwk-dlamj}
are finite sums of exponentials, so the score and observed Fisher information
$I_{\mathrm{obs}}(\hat\blam) = -\nabla^2 \ell(\hat\blam)$
are analytically tractable.  Standard errors are obtained as
$\widehat{\mathrm{se}}(\hat\lambda_j)
= \bigl[I_{\mathrm{obs}}^{-1}(\hat\blam)\bigr]_{jj}^{1/2}$.

In practice, we maximize~$\ell(\blam)$ using the L-BFGS-B algorithm
with box constraints $\lambda_j > 0$ for $j = 1,\ldots,m$.  As a
fallback for robustness, we re-optimize with the
Nelder--Mead simplex method on the reparametrized scale
$\log\lambda_j \in \Rset$, which enforces positivity without box
constraints; the solution yielding the higher log-likelihood is
retained.

\paragraph{Starting values.}
For a parallel system of $m$ i.i.d.\ $\mathrm{Exp}(\lambda)$ components,
$\E[T] = H_m / \lambda$ where $H_m = \sum_{k=1}^{m} 1/k$ is the $m$-th
harmonic number.  Writing $\hat\lambda_0 = H_m / \bar{t}$ for the
method-of-moments estimate (where $\bar{t}$ is the sample mean of the
observed system lifetimes), we initialize with a spread:
\begin{equation}\label{eq:spread-init}
  \lambda_j^{(0)} = \hat\lambda_0 \cdot
  \Bigl(\tfrac{1}{2} + \tfrac{j-1}{m-1}\Bigr), \qquad j = 1,\ldots,m.
\end{equation}
This places the initial rates in the correct order of magnitude while
breaking the permutation symmetry of the Scheme~0 log-likelihood.
A uniform initialization $\lambda_j^{(0)} = \hat\lambda_0$ for all~$j$
sits at a saddle point of the likelihood surface---where all rates are
equal---and gradient-based optimizers fail to escape.

% ===================================================================
\subsection{EM algorithm}
\label{sub:em-algorithm}
% ===================================================================

The censoring equivalence (\Cref{prop:censoring-equiv})
suggests a natural data-augmentation scheme: augment each observed
system lifetime~$T$ with the latent identity $J = \argmax_j T_j$ and
the unobserved component lifetimes $T_1,\ldots,T_m$.  This leads to an
EM algorithm in which both the E-step and the M-step admit closed-form
expressions.

\begin{proposition}[EM algorithm for exponential parallel systems]
  \label{prop:em-algorithm}
  Let $t_1,\ldots,t_n$ be $n$ i.i.d.\ observations from
  $T = \max(T_1,\ldots,T_m)$ with $T_j \sim \mathrm{Exp}(\lambda_j)$,
  $\lambda_j > 0$.  Define the complete data as
  $(T_1,\ldots,T_m,\, J)$ where $J = \argmax_j T_j$.

  Given current parameter estimates $\blam^{(s)}$, the EM update
  $\blam^{(s)} \mapsto \blam^{(s+1)}$ proceeds as follows.

  \medskip\noindent
  \textbf{E-step.}\enspace
  For each observation $t_i$ and each component $j = 1,\ldots,m$,
  compute:
  \begin{enumerate}
    \item[\textup{(i)}]
      The posterior probability that component $j$ was the last to fail:
      \begin{equation}
        \label{eq:em-posterior}
        p_{ij}
        \;\coloneqq\; P\bigl(J = j \mid T = t_i;\, \blam^{(s)}\bigr)
        = \frac{w_j\bigl(t_i;\, \blam^{(s)}\bigr)}
              {f_T\bigl(t_i;\, \blam^{(s)}\bigr)}\,.
      \end{equation}

    \item[\textup{(ii)}]
      The conditional mean of $T_j$ given that component $j$ was the
      last to fail:
      \begin{equation}
        \label{eq:em-exact}
        \E\bigl[T_j \mid T = t_i,\, J = j;\, \blam^{(s)}\bigr] = t_i\,.
      \end{equation}

    \item[\textup{(iii)}]
      The conditional mean of $T_j$ given that component $j$ was
      \emph{not} the last to fail:
      \begin{equation}
        \label{eq:em-truncated}
        \E\bigl[T_j \mid T = t_i,\, J \neq j;\, \blam^{(s)}\bigr]
        = \E\bigl[T_j \mid T_j < t_i;\, \lambda_j^{(s)}\bigr]
        = \frac{1}{\lambda_j^{(s)}}
          - \frac{t_i}{e^{\lambda_j^{(s)}\, t_i} - 1}\,.
      \end{equation}
  \end{enumerate}
  The overall conditional expectation of $T_j$ given observation $t_i$
  is then
  \begin{equation}
    \label{eq:em-etj}
    \hat\tau_{ij}
    \;\coloneqq\; \E\bigl[T_j \mid T = t_i;\, \blam^{(s)}\bigr]
    = p_{ij}\cdot t_i
      + (1 - p_{ij})\biggl[
        \frac{1}{\lambda_j^{(s)}}
        - \frac{t_i}{e^{\lambda_j^{(s)}\, t_i} - 1}
      \biggr].
  \end{equation}

  \medskip\noindent
  \textbf{M-step.}\enspace
  Update each rate as
  \begin{equation}
    \label{eq:em-mstep}
    \lambda_j^{(s+1)}
    = \frac{n}{\displaystyle\sum_{i=1}^{n} \hat\tau_{ij}}\,,
    \qquad j = 1,\ldots,m.
  \end{equation}
\end{proposition}

\begin{proof}
We verify each component of the E-step and derive the M-step.

\smallskip\noindent
\emph{Part~\textup{(i)}.}\enspace
By \Cref{lem:density-decomp},
$w_j(t;\blam)/f_T(t;\blam) = P(J = j \mid T = t;\blam)$.
Evaluating at $\blam = \blam^{(s)}$ gives~\eqref{eq:em-posterior}.

\smallskip\noindent
\emph{Part~\textup{(ii)}.}\enspace
Conditioning on $J = j$ means that component~$j$ achieved the maximum:
$T_j = T = t_i$.  Hence
$\E[T_j \mid T = t_i,\, J = j] = t_i$.

\smallskip\noindent
\emph{Part~\textup{(iii)}.}\enspace
Conditioning on $J \neq j$ means $T_j < T = t_i$, so $T_j$ follows
an $\mathrm{Exp}(\lambda_j)$ distribution truncated above at~$t_i$.
Its conditional density on $(0,\, t_i)$ is
$g(u) = \lambda_j\, e^{-\lambda_j u}\big/\bigl(1 - e^{-\lambda_j t_i}\bigr)$.
The conditional mean is
\begin{align*}
  \E[T_j \mid T_j < t_i]
  &= \frac{\displaystyle\int_0^{t_i} u\,\lambda_j\, e^{-\lambda_j u}\,du}
          {1 - e^{-\lambda_j t_i}}
  = \frac{1/\lambda_j - (1/\lambda_j + t_i)\,e^{-\lambda_j t_i}}
          {1 - e^{-\lambda_j t_i}}
  = \frac{1}{\lambda_j}
    - \frac{t_i}{e^{\lambda_j t_i} - 1}\,,
\end{align*}
where the numerator follows from integration by parts and the final
form is obtained by dividing numerator and denominator by
$e^{-\lambda_j t_i}$.

\smallskip\noindent
\emph{M-step derivation.}\enspace
The complete-data log-likelihood for a single system~$i$ with augmented
data $(T_{i1},\ldots,T_{im},\, J_i)$ is
\[
  \ell_c^{(i)}(\blam)
  = \sum_{j=1}^{m}\bigl(\ln \lambda_j - \lambda_j\, T_{ij}\bigr)
    + \text{terms not depending on~$\blam$},
\]
since each $T_{ij} \sim \mathrm{Exp}(\lambda_j)$ independently.  The
complete-data sufficient statistic for~$\lambda_j$ is
$\sum_{i=1}^n T_{ij}$, and the complete-data MLE is
$\hat\lambda_j = n \big/ \sum_{i} T_{ij}$.  Taking the conditional
expectation given the observed data and current parameters yields
$\E\bigl[\sum_i T_{ij} \mid t_1,\ldots,t_n;\, \blam^{(s)}\bigr]
= \sum_i \hat\tau_{ij}$.  The M-step maximizes the expected
complete-data log-likelihood, giving~\eqref{eq:em-mstep}.

\smallskip\noindent
\emph{Monotonicity.}\enspace
By the general theory of \citet{DempsterLairdRubin1977},
$\ell\bigl(\blam^{(s+1)}\bigr) \geq \ell\bigl(\blam^{(s)}\bigr)$ at
every iteration, with equality if and only if $\blam^{(s)}$ is a fixed
point of the EM map.  The complete-data model belongs to the exponential
family, so each M-step has a unique closed-form maximum.
\end{proof}

\begin{remark}[Closed-form character]
  \label{rem:em-closed-form}
  Both the E-step and M-step require only elementary operations:
  the E-step evaluates $w_j$ via the inclusion-exclusion expansion
  ($2^{m-1}$ terms per component) and the truncated exponential
  mean~\eqref{eq:em-truncated}; the M-step computes a reciprocal
  of a weighted average.  No numerical integration or Monte Carlo
  sampling is needed at any step.  This contrasts with the stochastic
  EM of \citet{NgNavarroBalakrishnan2017}, which requires simulation
  in the E-step for general system signatures, and with the Bayesian
  approach of \citet{BheringPereiraPolpo2014}, which uses MCMC.
\end{remark}

\begin{algorithm}[t]
  \caption{EM algorithm for exponential parallel system MLE}
  \label{alg:em}
  \begin{algorithmic}[1]
    \Require Observations $t_1,\ldots,t_n > 0$;\;
             number of components~$m$;\;
             tolerance $\varepsilon > 0$
    \Ensure  MLE $\hat\blam = (\hat\lambda_1,\ldots,\hat\lambda_m)$
    \State $\hat\lambda_0 \gets H_m / \bar{t}$
      \Comment{$H_m = \sum_{k=1}^m 1/k$,\; $\bar{t} = n^{-1}\sum_i t_i$}
    \State Initialize $\lambda_j^{(0)} \gets \hat\lambda_0 \cdot
      \bigl(\tfrac{1}{2} + \tfrac{j-1}{m-1}\bigr)$ for $j = 1,\ldots,m$
      \Comment{spread init~\eqref{eq:spread-init}}
    \State $s \gets 0$
    \Repeat
      \For{$i = 1,\ldots,n$;\; $j = 1,\ldots,m$}
        \Comment{\textbf{E-step}}
        \State $p_{ij} \gets w_j(t_i;\blam^{(s)}) \,\big/\, f_T(t_i;\blam^{(s)})$
        \State $\mu_{ij} \gets 1/\lambda_j^{(s)} - t_i\big/(e^{\lambda_j^{(s)} t_i} - 1)$
          \Comment{truncated mean}
        \State $\hat\tau_{ij} \gets p_{ij}\cdot t_i + (1 - p_{ij})\cdot \mu_{ij}$
      \EndFor
      \For{$j = 1,\ldots,m$}
        \Comment{\textbf{M-step}}
        \State $\lambda_j^{(s+1)} \gets n \,\big/\, \sum_{i=1}^n \hat\tau_{ij}$
      \EndFor
      \State $s \gets s + 1$
    \Until{$\bigl|\ell(\blam^{(s)}) - \ell(\blam^{(s-1)})\bigr| \big/ \bigl|\ell(\blam^{(s-1)})\bigr| < \varepsilon$}
    \State \Return $\hat\blam \gets \blam^{(s)}$
  \end{algorithmic}
\end{algorithm}

% ===================================================================
\subsection{Practical considerations}
\label{sub:practical}
% ===================================================================

\paragraph{Convergence criterion.}
We terminate the EM when the relative change in log-likelihood falls
below a threshold~$\varepsilon$:
$|\ell(\blam^{(s)}) - \ell(\blam^{(s-1)})| / |\ell(\blam^{(s-1)})| < \varepsilon$.
In our simulations (\Cref{sec:simulations}), $\varepsilon = 10^{-8}$
was used, and convergence typically required 50--200 iterations.

\paragraph{Agreement of methods.}
Both the direct maximization and the EM
algorithm target the same log-likelihood~\eqref{eq:loglik-exact}.  Since
\Cref{thm:identifiability} guarantees a unique global maximum
(up to label permutation), they converge to the same point estimate.
The direct method supplies the observed Fisher information at the MLE
as a byproduct; when using the EM, standard errors are obtained by
evaluating the Hessian of~$\ell$ at the converged estimate~$\hat\blam$.

\paragraph{Computational cost.}
Each evaluation of $f_T(t;\blam)$ requires
the inclusion-exclusion expansion for each of $m$ components, involving
$2^{m-1}$ terms per component.  Over $n$ observations, a single
likelihood evaluation costs $O(n \cdot m \cdot 2^{m-1})$.  The
exponential scaling in~$m$ is the binding constraint: for the
exponential model, systems with $m \leq 15$--$20$ are computationally
feasible on standard hardware.

% ===================================================================
\subsection{Extension to censored system data}
\label{sub:censored-extension}
% ===================================================================

In practice, the system lifetime~$T$ may itself be subject to censoring.
The inclusion-exclusion machinery of \Cref{sec:likelihood}
accommodates all standard censoring types in closed form.

If system~$i$ is known only to survive past time~$\tau_i$
(i.e., $T_i > \tau_i$), its likelihood contribution is the system
survival function
$S_T(\tau_i;\blam) = 1 - F_T(\tau_i;\blam)$,
available in closed form via \Cref{lem:ie}.
If the system failure is known to lie in an interval $(a_i,\, b_i]$,
the likelihood contribution is
$F_T(b_i;\blam) - F_T(a_i;\blam)$,
a difference of two inclusion-exclusion expansions.

For a dataset containing exact, right-censored, and interval-censored
observations, the mixed log-likelihood~\eqref{eq:loglik}
applies directly, with every term in closed form.  Both the direct
maximization and EM approaches extend to this setting without
modification.  For the EM, the E-step conditions on the appropriate
censoring information: exact observations contribute as in
\Cref{prop:em-algorithm}; censored observations contribute the
conditional expectation of the component lifetimes given the
censoring constraint, which involves only exponential and truncated
exponential calculations.

% Section 5: Monte Carlo Study
\section{Monte Carlo Study}
\label{sec:simulations}

We evaluate the finite-sample properties of the direct MLE developed in
\Cref{sec:estimation} through three Monte Carlo experiments targeting
sample size effects, system-level censoring, and the number of system
components.  The central empirical finding is an information asymmetry
across components: the precision of $\hat{\lambda}_j$ is governed by
$P(T_j = T)$, the probability that component~$j$ is the last to fail.

%% ===================================================================
\subsection{Simulation design}
\label{sec:sim-design}

\paragraph{Data generation.}
For each replication, we generate $n$ independent system lifetimes from
$T = \max(T_1, \ldots, T_m)$ with $T_j \sim \mathrm{Exp}(\lambda_j)$
under Scheme~0 (black-box observation): only the system lifetime~$T$ is
recorded, with no component-level information.
The system lifetime is then processed through an observation mechanism
(exact or right-censored; see below).

\paragraph{Estimation.}
The log-likelihood~\eqref{eq:loglik} is maximized using L-BFGS-B
with positivity constraints on $\blam$, with Nelder--Mead fallback on a
log-transformed parameterization.  Since the Scheme~0 likelihood is
permutation-symmetric (\Cref{thm:identifiability}), starting values
are spread asymmetrically around a method-of-moments center to break
the symmetry and avoid the equal-rate saddle point.  Standard errors
are computed from the observed Fisher information via numerical
differentiation.

\paragraph{Metrics.}
For each configuration, we run $R = 500$ Monte Carlo replications with
base seed $2026$ and report the following over converged replications:
\begin{itemize}
  \item \emph{Bias}: $R^{-1}\sum_{r} (\hat{\lambda}_j^{(r)} - \lambda_j)$;
  \item \emph{RMSE}: $\bigl[R^{-1}\sum_{r}(\hat{\lambda}_j^{(r)} - \lambda_j)^2\bigr]^{1/2}$;
  \item \emph{Coverage}: proportion of replications in which $\lambda_j$
    falls within the 95\% Wald interval
    $\hat{\lambda}_j \pm 1.96\,\widehat{\mathrm{se}}_j$;
  \item \emph{SE ratio}: mean estimated standard error divided by the
    empirical standard deviation of $\hat{\lambda}_j$ across replications
    (values near~1 indicate well-calibrated uncertainty quantification).
\end{itemize}

\paragraph{Baseline configuration.}
Unless otherwise noted, we use $m = 3$ components with rates
$\blam = (0.2, 0.3, 0.5)$ (listed in ascending order throughout),
sample size $n = 200$, and right-censoring
threshold $\tau = 15$ (yielding negligible censoring at approximately~6\%).
All simulations are implemented in R.

\paragraph{Label matching.}
Since the MLE identifies rates only as a multiset
(\Cref{thm:identifiability}), we sort each replicate's estimated rates
in ascending order and compare them to the sorted true rates.
This is necessary because the Scheme~0 likelihood is
permutation-symmetric: the optimizer may recover the correct
multiset in any order.


%% ===================================================================
\subsection{Experiment 1: Effect of sample size}
\label{sec:sim-sample-size}

We vary $n \in \{50, 100, 200, 500\}$ while holding all other parameters
at their baseline values.  \Cref{tab:sample-size} reports the results;
\Cref{fig:rmse-sample-size} displays the RMSE trajectories.

\begin{table}[t]
\centering
\caption{%
  MLE performance as a function of sample size ($m = 3$,
  $\blam = (0.2, 0.3, 0.5)$, $\tau = 15$,
  $R = 500$ replications).
}
\label{tab:sample-size}
\small
\begin{tabular}{rl rrrr}
\toprule
$n$ & Parameter & Bias & RMSE & Cov.\ (\%) & SE ratio \\
\midrule
  50 & $\lambda_1 = 0.2$ & $+$0.024  & 0.056 & 96.2 & 2.53 \\
     & $\lambda_2 = 0.3$ & $+$0.040  & 0.118 & 100.0 & 2.85 \\
     & $\lambda_3 = 0.5$ & $+$0.489  & 2.277 & 94.4 & 0.71 \\
\addlinespace
 100 & $\lambda_1 = 0.2$ & $+$0.015  & 0.044 & 94.5 & 2.36 \\
     & $\lambda_2 = 0.3$ & $+$0.033  & 0.097 & 100.0 & 3.00 \\
     & $\lambda_3 = 0.5$ & $+$0.203  & 1.197 & 95.4 & 0.67 \\
\addlinespace
 200 & $\lambda_1 = 0.2$ & $+$0.014  & 0.037 & 95.2 & 2.46 \\
     & $\lambda_2 = 0.3$ & $+$0.021  & 0.085 & 99.8 & 2.92 \\
     & $\lambda_3 = 0.5$ & $+$0.074  & 0.730 & 96.2 & 0.59 \\
\addlinespace
 500 & $\lambda_1 = 0.2$ & $+$0.010  & 0.027 & 92.8 & 2.31 \\
     & $\lambda_2 = 0.3$ & $+$0.009  & 0.072 & 99.0 & 2.54 \\
     & $\lambda_3 = 0.5$ & $+$0.008  & 0.159 & 99.8 & 1.48 \\
\bottomrule
\end{tabular}
\end{table}

\begin{figure}[t]
  \centering
  \includegraphics[width=0.78\textwidth]{../plots/sample_size/rmse.pdf}
  \caption{%
    RMSE of the MLE as a function of sample size for each component rate.
    The fastest component ($\lambda_3 = 0.5$) exhibits dramatically
    elevated RMSE at $n = 50$, driven by heavy right tails in the sampling
    distribution.  By $n = 200$, all three components achieve moderate RMSE,
    and the $\sqrt{n}$ convergence rate is visible.
  }
  \label{fig:rmse-sample-size}
\end{figure}

The MLE is consistent across all three components: bias diminishes toward
zero as $n$ grows, and RMSE decreases at a rate consistent with $n^{-1/2}$.
Coverage of the 95\% Wald interval is within simulation error of the
nominal level for $n \geq 100$.

The most striking feature is the behavior of $\hat{\lambda}_3$, the
fastest component.  At $n = 50$, its RMSE is~2.28---more than
$4.5$ times the true parameter value---while $\hat{\lambda}_1$ and
$\hat{\lambda}_2$ have RMSE of~0.056 and~0.118 respectively.  This
order-of-magnitude disparity reflects the information asymmetry inherent in
parallel systems: $\lambda_3 = 0.5$ is the fastest component, meaning it
rarely determines the system lifetime ($P(T_3 = T) = 0.161$ for this
configuration).  Its failure time~$T_3$ is deeply left-censored by the
system observation $T = \max T_j$, so the data carry little direct
information about~$\lambda_3$.

Despite the elevated RMSE at $n = 50$, Wald coverage remains
near~95\% or above for all parameters.  The SE ratio column reveals an
important nuance of sorted-matching inference: the Hessian-based
standard errors are computed for the \emph{unsorted} MLE, but we report
performance against \emph{sorted} estimates.  The sorted middle
parameter ($\hat\lambda_2$) has less sampling variability than an
unsorted parameter (it is sandwiched between the extremes), producing
SE ratios above~1 and overcoverage.  The sorted extreme
($\hat\lambda_3$) has SE ratio below~1, indicating that the
Hessian underestimates the tail variability of the largest order
statistic.  As $n$ grows, the sorting effect diminishes: at $n = 500$,
all SE ratios are between 1.5 and 2.5.
Profile likelihood or bootstrap intervals would eliminate this artifact;
we leave their implementation to future work.


%% ===================================================================
\subsection{Experiment 2: Effect of system-level censoring}
\label{sec:sim-censoring}

Holding $m = 3$, $n = 200$, and $\blam = (0.2, 0.3, 0.5)$, we consider
three observation schemes:
(i)~exact, with no censoring ($\tau = \infty$);
(ii)~right-censored at $\tau = 8$ (${\approx}29\%$ censored); and
(iii)~right-censored at $\tau = 4$ (${\approx}67\%$ censored).
\Cref{tab:censoring} reports the results; \Cref{fig:rmse-censoring}
displays the RMSE comparison.

\begin{table}[t]
\centering
\caption{%
  MLE performance under different observation schemes ($m = 3$,
  $\blam = (0.2, 0.3, 0.5)$, $n = 200$, $R = 500$).
}
\label{tab:censoring}
\small
\begin{tabular}{ll rrrr}
\toprule
Scheme & Parameter & Bias & RMSE & Cov.\ (\%) & Conv.\ (\%) \\
\midrule
Exact
  & $\lambda_1 = 0.2$ & $+$0.013 & 0.034 & 95.6 & 100 \\
  & $\lambda_2 = 0.3$ & $+$0.017 & 0.080 & 98.8 &  \\
  & $\lambda_3 = 0.5$ & $+$0.081 & 0.702 & 96.4 &  \\
\addlinespace
Right, $\tau\!=\!8$
  & $\lambda_1 = 0.2$ & $+$0.018 & 0.049 & 91.7 & 100 \\
  & $\lambda_2 = 0.3$ & $+$0.032 & 0.100 & 99.8 &  \\
  & $\lambda_3 = 0.5$ & $+$0.066 & 0.747 & 97.2 &  \\
\addlinespace
Right, $\tau\!=\!4$
  & $\lambda_1 = 0.2$ & $+$0.027 & 0.072 & 89.5 & 99.6 \\
  & $\lambda_2 = 0.3$ & $+$0.050 & 0.145 & 100.0 &  \\
  & $\lambda_3 = 0.5$ & $+$0.254 & 1.916 & 100.0 &  \\
\bottomrule
\end{tabular}
\end{table}

\begin{figure}[t]
  \centering
  \includegraphics[width=0.78\textwidth]{../plots/censoring/rmse.pdf}
  \caption{%
    RMSE across observation schemes.  Exact observation and mild
    right-censoring ($\tau = 8$) yield low RMSE.  Aggressive
    right-censoring ($\tau = 4$) degrades all estimates substantially.
  }
  \label{fig:rmse-censoring}
\end{figure}

Mild right-censoring ($\tau = 8$, ${\approx}29\%$ censored) has a
modest impact: RMSE increases by 6--44\% relative to exact
observation, and coverage remains acceptable.  Aggressive censoring
($\tau = 4$, ${\approx}67\%$ censored) substantially degrades
estimation: for the fastest component ($\lambda_3 = 0.5$), RMSE nearly
triples from~0.70 to~1.92.  All three components exhibit positive
bias, as right-censoring truncates the informative upper tail of the
system lifetime distribution.  Coverage for~$\lambda_1$ drops to 89.5\%
under aggressive censoring, while the middle and fastest components
maintain $\geq 97\%$ coverage (due to the conservative SE estimates
discussed above).

Interval-censored system data (periodic inspection) is also
accommodated by the likelihood framework of \Cref{sec:likelihood},
but we defer its simulation study to future work, as the direct
optimizer exhibited convergence difficulties in preliminary
experiments.  The EM algorithm of \Cref{prop:em-algorithm} may be
better suited to this setting.


%% ===================================================================
\subsection{Experiment 3: Effect of the number of components}
\label{sec:sim-components}

We vary the number of components: $m = 2$ with $\blam = (0.3, 0.5)$,
$m = 3$ with $\blam = (0.2, 0.3, 0.5)$ (the baseline), and $m = 5$
with $\blam = (0.1, 0.2, 0.3, 0.4, 0.5)$.  All use $n = 200$ and
$\tau = 15$.

This experiment reveals the paper's central empirical finding: the
\emph{information asymmetry} across components in a parallel system.

\paragraph{The governing quantity: $P(T_j = T)$.}
The probability that component~$j$ is the last to fail,
\begin{equation}
\label{eq:prob-last-sim}
  P(T_j = T)
    = \lambda_j \sum_{S \subseteq \{1,\ldots,m\}\setminus\{j\}}
      \frac{(-1)^{|S|}}{\lambda_j + \lambda(S)}\,,
\end{equation}
governs how much system-level information is available
about component~$j$: when $T_j = T$, component~$j$'s failure time is
observed exactly; when $T_j < T$, it is left-censored at the system
lifetime.  The conditional distribution of component lifetimes given
system failure---the ``mean past lifetime'' studied by
\citet{Asadi2006}---provides the distributional perspective on this
information asymmetry.  \Cref{tab:prob-last} reports exact values for each
configuration.

\begin{table}[t]
\centering
\caption{%
  Probability $P(T_j = T)$ that component~$j$ determines the system
  lifetime, computed exactly via~\eqref{eq:prob-last}.
  Slow components (small~$\lambda_j$) dominate; fast components rarely do.
  Parameters listed in ascending order to match the sorted-estimate
  convention used in all simulation tables.
}
\label{tab:prob-last}
\small
\begin{tabular}{l ccccc}
\toprule
& $\lambda_1$ & $\lambda_2$ & $\lambda_3$ & $\lambda_4$ & $\lambda_5$ \\
\midrule
$m = 2$: $(0.3,\, 0.5)$
  & 0.625 & 0.375 & --- & --- & --- \\
$m = 3$: $(0.2,\, 0.3,\, 0.5)$
  & 0.514 & 0.325 & 0.161 & --- & --- \\
$m = 5$: $(0.1,\, 0.2,\, 0.3,\, 0.4,\, 0.5)$
  & 0.536 & 0.229 & 0.120 & 0.070 & 0.044 \\
\bottomrule
\end{tabular}
\end{table}

The asymmetry is stark.  In the $m = 5$ system, the slowest component
($\lambda_1 = 0.1$) is the last to fail in 53.6\% of realizations, while
the fastest ($\lambda_5 = 0.5$) is last in only 4.4\%---a ratio of
roughly $12{:}1$.

\paragraph{Five-component results.}
\Cref{tab:components-m5} reports the estimation results for $m = 5$
(417 of 500 replications converged).
\Cref{fig:sampling-m5} displays the sampling distributions.

\begin{table}[t]
\centering
\caption{%
  MLE performance for the $m = 5$ parallel system
  ($\blam = (0.1, 0.2, 0.3, 0.4, 0.5)$, $n = 200$, $R = 500$,
  417 converged).
  RMSE varies by nearly three orders of magnitude across components,
  closely tracking $P(T_j = T)$.
}
\label{tab:components-m5}
\small
\begin{tabular}{l ccccc}
\toprule
Parameter & True & $P(T_j\!=\!T)$ & Bias & RMSE & Cov.\ (\%) \\
\midrule
$\lambda_1$ & 0.10 & 0.536 & $+$0.009 & 0.026 & 92.8 \\
$\lambda_2$ & 0.20 & 0.229 & $+$0.037 & 0.087 & 100.0 \\
$\lambda_3$ & 0.30 & 0.120 & $+$0.011 & 0.094 & 99.5 \\
$\lambda_4$ & 0.40 & 0.070 & $+$0.046 & 0.271 & 100.0 \\
$\lambda_5$ & 0.50 & 0.044 & $+$1.028 & 7.419 & 100.0 \\
\bottomrule
\end{tabular}
\end{table}

\begin{figure}[t]
  \centering
  \includegraphics[width=\textwidth]{../plots/sampling_dists/components_m5_sampling_dists.pdf}
  \caption{%
    Sampling distributions of $\hat{\blam}$ for the $m = 5$ parallel
    system ($n = 200$).  The panels differ
    dramatically in horizontal scale.  The slowest component
    ($\hat{\lambda}_1$) has a tight, approximately Gaussian distribution.
    The fastest components ($\hat{\lambda}_4$, $\hat{\lambda}_5$) exhibit
    extreme right skew, reflecting near-total left-censoring of these
    components' failure times.
  }
  \label{fig:sampling-m5}
\end{figure}

For $\lambda_1 = 0.1$ (the slowest component, $P(T_1 = T) = 0.536$),
the MLE achieves RMSE of~0.026, bias of~$+0.009$, and coverage
of~92.8\%---near-ideal behavior.  For
$\lambda_5 = 0.5$ (the fastest, $P(T_5 = T) = 0.044$), the RMSE rises
to~7.42, over $14.8\times$ the true parameter value, and the sampling
distribution (\Cref{fig:sampling-m5}) exhibits extreme right skew.

The RMSE ratio between the worst- and best-estimated components is
$7.42 / 0.026 \approx 285$, spanning nearly three orders of magnitude.
Coverage remains at or above~93\% for all parameters, though
the sorted-matching effect inflates coverage for the interior components
($\lambda_2$--$\lambda_4$ all show 99.5--100\%).  The information
asymmetry under pure black-box observation is substantially starker
than would be observed with partial component identity information.

\paragraph{Comparison across $m$.}
For $m = 2$ ($\blam = (0.3, 0.5)$), the information asymmetry is mild:
$P(T_j = T)$ values of~0.625 and~0.375 yield RMSE of~0.044 and~0.118,
with coverage of~96.2\% for both parameters.  The three-component
baseline ($m = 3$) shows the onset of asymmetry, with $\hat{\lambda}_3$
having RMSE roughly $20\times$ that of $\hat{\lambda}_1$ (0.730 versus
0.037).  The $m = 5$ case makes the asymmetry dramatic and
unmistakable.


%% ===================================================================
\subsection{EM algorithm comparison}
\label{sec:sim-em}

To validate the EM algorithm of \Cref{prop:em-algorithm}, we compare
it against the direct MLE on the baseline configuration ($m = 3$,
$n = 200$, $\blam = (0.2, 0.3, 0.5)$, Scheme~0).  Over 20~random
seeds, the maximum log-likelihood discrepancy between the two methods
is~$1.3 \times 10^{-3}$, and the maximum parameter discrepancy
(sorted) is~$2.7 \times 10^{-2}$, confirming that both converge to the
same global maximum.  The EM requires a median of~200 iterations
(range 95--490) with tolerance $\varepsilon = 10^{-8}$ and is
approximately $2\times$ slower than direct L-BFGS-B maximization for
this problem size.  For $m = 5$, both methods again agree to within
log-likelihood discrepancies of~$3 \times 10^{-3}$, though the EM
requires 250--1660 iterations.


%% ===================================================================
\subsection{Information efficiency}
\label{sec:sim-fisher}

The information asymmetry documented in the preceding experiments can be
quantified through the Fisher information.  For complete component data
(Scheme~2: all $T_1, \ldots, T_m$ observed), the Fisher information
for $\lambda_j$ from $n$ independent observations is
\begin{equation}\label{eq:fisher-complete}
  I_j^{\mathrm{complete}} = n / \lambda_j^2,
\end{equation}
the standard result for exponential MLEs.  Under Scheme~0, the observed
Fisher information at the MLE---computed as $\hat{I}(\hat\blam) =
-H(\hat\blam)$, where $H$ is the Hessian of the log-likelihood---is
necessarily smaller (in a matrix sense), since system-level data carry
less information than component-level data.

We define the \emph{per-component information efficiency} as the ratio of
diagonal entries:
\begin{equation}\label{eq:efficiency}
  e_j = \hat{I}_{jj}(\hat\blam) \big/ I_j^{\mathrm{complete}},
\end{equation}
and the \emph{joint efficiency} as the determinant ratio
$\det\bigl(\hat{I}(\hat\blam)\bigr) \big/ \det\bigl(I^{\mathrm{complete}}\bigr)$.
Values near~1 indicate that system-level observation captures nearly all
component-level information; values near~0 indicate severe information
loss.

\Cref{tab:fisher-info} reports median efficiencies over $R = 200$
replications for $n = 200$ observations, computed at the MLE with sorted
parameter matching.  \Cref{fig:fisher-efficiency} displays the
relationship between $P(T_j = T)$ and the per-component efficiency.

\begin{table}[t]
\centering
\caption{%
  Information efficiency of Scheme~0 relative to complete data
  ($n = 200$, $R = 200$ replications).  The per-component
  efficiency~$e_j$ closely tracks $P(T_j = T)$.  The joint efficiency
  (determinant ratio) collapses rapidly as~$m$ grows.
}
\label{tab:fisher-info}
\small
\begin{tabular}{rl ccc}
\toprule
$m$ & Parameter & $P(T_j\!=\!T)$ & Median $e_j$ & Mean $e_j$ \\
\midrule
2 & $\lambda_1 = 0.3$ & 0.625 & 0.79 & 0.76 \\
  & $\lambda_2 = 0.5$ & 0.375 & 0.19 & 0.33 \\
  & \multicolumn{3}{l}{\textit{Joint efficiency (det ratio):
      $5.8 \times 10^{-2}$}} \\
\addlinespace
3 & $\lambda_1 = 0.2$ & 0.514 & 0.56 & 0.59 \\
  & $\lambda_2 = 0.3$ & 0.325 & 0.16 & 0.33 \\
  & $\lambda_3 = 0.5$ & 0.161 & 0.11 & 0.18 \\
  & \multicolumn{3}{l}{\textit{Joint efficiency: $3.9 \times 10^{-4}$}} \\
\addlinespace
5 & $\lambda_1 = 0.1$ & 0.536 & 0.58 & 0.53 \\
  & $\lambda_2 = 0.2$ & 0.229 & 0.05 & 0.19 \\
  & $\lambda_3 = 0.3$ & 0.120 & 0.05 & 0.14 \\
  & $\lambda_4 = 0.4$ & 0.070 & 0.06 & 0.08 \\
  & $\lambda_5 = 0.5$ & 0.044 & 0.03 & 0.06 \\
  & \multicolumn{3}{l}{\textit{Joint efficiency: $6.1 \times 10^{-10}$}} \\
\bottomrule
\end{tabular}
\end{table}

\begin{figure}[t]
  \centering
  \includegraphics[width=0.78\textwidth]{../plots/fisher_info/efficiency.pdf}
  \caption{%
    Per-component information efficiency $e_j$ (mean $\pm$ 1~SD) versus
    $P(T_j = T)$ for $m = 2, 3, 5$ parallel systems ($n = 200$).
    The slowest component (highest $P(T_j = T)$) retains 53--79\% of
    the complete-data information; the fastest retains only 3--19\%.
    The near-linear relationship confirms that $P(T_j = T)$ governs
    the information content of system-level observation.
  }
  \label{fig:fisher-efficiency}
\end{figure}

Three findings emerge.  First, the per-component efficiency~$e_j$ is
approximately proportional to~$P(T_j = T)$: the slowest component
retains 53--79\% of complete-data information, while the fastest retains
only 3--19\%.  Second, the joint efficiency collapses exponentially
with~$m$: from $5.8 \times 10^{-2}$ at $m = 2$ to $6.1 \times 10^{-10}$
at $m = 5$.  This reflects the compounding effect of information loss
across components---the system-level observation is a maximally lossy
compression of the full component data.  Third, the mean efficiency
exceeds the median for fast components, reflecting the heavy right tail
of the observed information in those cases (a consequence of the
rare-event dependence: the few replications where the fast component
happens to be last yield disproportionately large information).

These results make the $P(T_j = T)$ asymmetry from
\Cref{sec:sim-components} quantitatively precise: the information
loss under Scheme~0 is governed component-by-component by the
probability of being the last to fail.  The joint efficiency collapse
further motivates the development of more informative observation
schemes (\Cref{sec:discussion}).


%% ===================================================================
\subsection{Summary}
\label{sec:sim-summary}

The four experiments establish the following empirical findings:
\begin{enumerate}
  \item \textbf{Consistency and asymptotic normality.}
    The direct MLE is consistent with bias vanishing as~$n$ grows, RMSE
    decreasing at the $\sqrt{n}$ rate, and 95\% Wald coverage at or
    above the nominal level for all sample sizes.  SE ratios
    deviate from~1 due to the sorted-matching convention (see below),
    but converge toward~1 as~$n$ grows.

  \item \textbf{Censoring effects are asymmetric.}
    Right-censoring degrades estimation quality in proportion to the
    censoring fraction, with positive bias and reduced coverage.
    Interval-censored system data, under direct maximization, exhibits
    convergence difficulties that require further investigation---potentially
    via the EM algorithm.

  \item \textbf{Information asymmetry governed by $P(T_j = T)$.}
    The probability that component~$j$ determines the system lifetime is
    the governing quantity for estimation precision.  In a five-component
    system, RMSE varies by nearly three orders of magnitude: the slowest
    component ($P(T_1 = T) \approx 0.54$) achieves RMSE of~0.026, while
    the fastest ($P(T_5 = T) \approx 0.04$) has RMSE of~7.42,
    effectively unidentifiable from $n = 200$ black-box observations.

  \item \textbf{Fisher information efficiency.}
    The per-component efficiency $e_j$ (ratio of Scheme~0 observed
    information to complete-data information) is approximately
    proportional to $P(T_j = T)$, ranging from 53--79\% for the
    slowest component to 3--19\% for the fastest.  The joint
    efficiency (determinant ratio) collapses from $5.8 \times 10^{-2}$
    at $m = 2$ to $6.1 \times 10^{-10}$ at $m = 5$.
\end{enumerate}

Finding~(3), quantified by~(4), carries a direct practical implication:
if precise estimates of all component rates are required, system-level
observation alone is insufficient for fast components.  Component-level
monitoring (Scheme~1 or higher) is needed.  This mirrors the dual
situation in series systems, where the weakest-link component dominates
and is well-estimated while stronger components are right-censored.  In
parallel systems, the strongest-link component dominates and is
well-estimated; weaker components are left-censored by the system
structure.

% Section 6: Discussion
\section{Discussion}
\label{sec:discussion}

\subsection{Summary}

We have shown that the incomplete-data structure of a parallel system
observation is component-level left-censoring with latent identity,
not masked cause of failure.  This censoring equivalence
(\Cref{prop:censoring-equiv}) provides the conceptual foundation for
a closed-form MLE via inclusion-exclusion, an EM algorithm with
analytical E- and M-steps, and an elementary identifiability proof for
exponential components.  The simulations of \Cref{sec:simulations}
confirm that the MLE is consistent with nominal Wald coverage, and
reveal a striking information asymmetry: the precision of
$\hat\lambda_j$ is governed by $P(T_j = T)$, the probability that
component~$j$ determines the system lifetime.

\subsection{The information asymmetry}

The quantity $P(T_j = T)$ emerges as the organizing principle for
parallel system estimation.  Slow components (small~$\lambda_j$) fail
last with high probability, so their failure times are observed nearly
exactly in most system observations.  Fast components (large~$\lambda_j$)
fail early and are deeply left-censored by the system structure, yielding
RMSE ratios exceeding~$150{:}1$ in the five-component case.

This has a direct practical implication.  If reliable estimates of all
component rates are needed, system-level observation alone is
insufficient for fast components.  Additional component-level
information---through periodic inspection (Scheme~1), partial
monitoring (Scheme~3), or degradation signatures---is required.  The
Fisher information analysis of \Cref{sec:sim-fisher} quantifies
this: the per-component efficiency under Scheme~0 is approximately
proportional to $P(T_j = T)$, and the joint efficiency (determinant
ratio) collapses exponentially with~$m$, from $5.8 \times 10^{-2}$ at
$m = 2$ to $6.1 \times 10^{-10}$ at $m = 5$.  Extending this
comparison to Scheme~1 (periodic inspection) and other intermediate
observation schemes is a natural direction for guiding monitoring
strategy.

The asymmetry mirrors the dual situation in series systems.  There,
the weakest-link component ($\argmin_j T_j$) dominates the system
lifetime and is well-estimated from system data, while stronger
components are right-censored.  In parallel systems, the
strongest-link component ($\argmax_j T_j$) dominates, and weaker
components are left-censored.  The censoring direction reverses, but
the structural pattern is the same: the component that determines the
system lifetime is the one about which the most is learned.

\subsection{Limitations}

\paragraph{Exponential assumption.}
The constant-hazard assumption of the exponential distribution is
restrictive for many applications.  The inclusion-exclusion expansion
and the closed-form integrals (\Cref{lem:ie,lem:integral-wj}) depend
on the exponential form of the component CDF.  Extension to Weibull
or gamma components would preserve the general structure of the
density decomposition (\Cref{lem:density-decomp}) and the EM
framework (\Cref{prop:em-algorithm}), but the closed-form integrals
would be replaced by numerical quadrature, and the identifiability
proof would require different tail analysis.

\paragraph{Computational scaling.}
The inclusion-exclusion expansion has $2^m$ terms, limiting practical
applicability to systems with roughly $m \leq 15$--$20$ components.
For larger systems, Monte Carlo estimation of the density, recursive
decomposition of the inclusion-exclusion sum, or restriction to i.i.d.\
component models would be needed.

\paragraph{Independence.}
The independence assumption underlies both the product-form CDF and the
identifiability proof (\Cref{thm:identifiability}).  This is not merely
a convenience: the \citet{BasuGhosh1980} identifiability result
\emph{requires} independence.  Under general dependence, the
distribution of $\max(T_1, \ldots, T_m)$ no longer uniquely determines
the marginal component distributions---the same system CDF can arise
from different combinations of marginals and copula.  This is the
parallel-system analogue of the \citet{Tsiatis1975} non-identifiability
for series systems, but it appears to be unstated in the literature:
recent work on identifiability under dependence
\citep{HiabuLoWilke2025} addresses the series case via exclusion
restrictions, while the parallel case under dependence remains open.
Copula extensions \citep{NavarroNgBalakrishnan2012} could model
dependence but would require either structural assumptions on the
copula family or additional observation beyond Scheme~0.

\paragraph{Interval-censored optimization.}
The simulation study (\Cref{sec:sim-censoring}) revealed convergence
difficulties for purely interval-censored system data under direct
maximization.  The EM algorithm may provide better numerical behavior
in this setting, but a systematic comparison was beyond the scope of
the present study.

\subsection{Extension to $k$-out-of-$n$ systems}
\label{sec:k-out-of-n}

The parallel system ($1$-out-of-$n$) is the extreme case of a
$k$-out-of-$n$ system: the system functions as long as at least $k$ of
its $n$ components survive.  Equivalently, the system fails when the
$(n - k + 1)$-th component fails:
\begin{equation}\label{eq:kofn-lifetime}
  T = T_{(n-k+1)},
\end{equation}
where $T_{(1)} \leq T_{(2)} \leq \cdots \leq T_{(n)}$ denote the order
statistics of the component lifetimes.  The parallel system has $k = 1$
(system lifetime $T = T_{(n)}$, the maximum), and the series system has
$k = n$ (system lifetime $T = T_{(1)}$, the minimum).

The censoring equivalence (\Cref{prop:censoring-equiv}) generalizes
naturally to this setting.

\begin{proposition}[$k$-out-of-$n$ censoring equivalence]
\label{prop:kofn-censoring}
Consider independent, continuous component lifetimes
$T_1, \ldots, T_n$ and a $k$-out-of-$n$ system with failure
time $T = T_{(n-k+1)}$.  Observing the system lifetime~$T$ under
Scheme~0 is informationally equivalent to the following component-level
censoring scheme:
\begin{enumerate}
  \item[\textup{(a)}] One component---the $(n-k+1)$-th to fail---is
    observed exactly: $T_J = T$ for a latent index~$J$.
  \item[\textup{(b)}] The $n - k$ components that failed before the system
    are left-censored at~$T$: for these indices~$j$, we know only that
    $T_j < T$.
  \item[\textup{(c)}] The $k - 1$ components still functioning at system
    failure are right-censored at~$T$: for these indices~$j$, we know
    only that $T_j > T$.
  \item[\textup{(d)}] The identities of the components in each group are
    latent.
\end{enumerate}
\end{proposition}

\begin{proof}
When the system fails at time $T = T_{(n-k+1)} = t$, exactly $n - k + 1$
components have failed (those with $T_j \leq t$) and exactly $k - 1$
components survive (those with $T_j > t$).  Among the failed components,
exactly one satisfies $T_j = t$ almost surely (by continuity of the
component distributions), and the remaining $n - k$ have $T_j < t$.
Under Scheme~0, the system-level observation~$T = t$ reveals only this
count structure, not which components belong to which group.  This yields
one exact, $n - k$ left-censored, and $k - 1$ right-censored component
observations with latent assignment.
\end{proof}

\Cref{prop:kofn-censoring} recovers \Cref{prop:censoring-equiv} when
$k = 1$: the parallel system has $n - 1$ left-censored and zero
right-censored observations.  At the other extreme, $k = n$ gives the
series system: zero left-censored and $n - 1$ right-censored
observations---exactly the competing risks structure underlying the
classical masked cause-of-failure problem.

\begin{remark}[Information hierarchy]
\label{rmk:info-hierarchy}
The proposition reveals a monotone information structure as $k$
varies.  Decreasing~$k$ (from series toward parallel) converts
right-censored components into left-censored ones, but the total number
of censored components remains $n - 1$ regardless of~$k$.  What changes
is the \emph{direction}: left-censoring (parallel) and right-censoring
(series) carry different amounts of information, giving rise to
the identifiability contrast between the two extremes
(\Cref{cor:duality}).  For intermediate~$k$, both censoring
directions coexist, and the information content interpolates.
\end{remark}

For the discrete analogue (geometric component lifetimes),
\citet{Jasinski2023} and \citet{Jasinski2024} have derived closed-form
MLEs for $k$-out-of-$n$ systems with homogeneous and heterogeneous
components, respectively.  \Cref{prop:kofn-censoring} provides the
continuous counterpart, establishing the censoring structure that
underlies estimation for any parametric family.

\subsection{Related developments}

The censoring framework developed here complements several recent
research directions.  \citet{Cramer2023} provides a comprehensive
review of ordered and censored lifetime data in reliability, situating
progressive censoring models (including the coherent system
representation of \citealt{CramerNavarro2015}) within a unified
framework.  \citet{QuNgMoon2024a,QuNgMoon2024b} develop nonparametric
tests for component distribution homogeneity from system data with known
signatures---a testing complement to our estimation approach.
\citet{BiswasGangulyMitra2025} propose semiparametric (piecewise-linear
cumulative hazard) models for load-sharing parallel systems, providing a
flexible alternative to the parametric exponential model studied here.
These directions---nonparametric testing, semiparametric models, and the
parametric censoring-based MLE developed in the present paper---represent
complementary approaches to the same fundamental problem of learning
about components from system-level data.

A structural connection also exists with the one-shot device testing
literature \citep{BalakrishnanSo2025}, where each test unit is
destroyed upon inspection and only pass/fail status at the inspection
time is recorded.  For a parallel system observed at a single time
point (current-status data), the Scheme~0 observation reduces to a
binary outcome---structurally identical to the one-shot device model
with multiple competing failure modes.  The two literatures have
developed largely independently; the censoring equivalence developed
here provides a bridge.

\subsection{Future work}

Several extensions arise naturally from the censoring framework
developed here.

\paragraph{Periodic inspection (Scheme~1).}
Under periodic inspection at times $\delta, 2\delta, 3\delta, \ldots$,
each component failure is localized to an inspection interval, yielding
component-level interval-censored data.  For a parallel system observed
under both Scheme~0 (system failure time~$T$) and Scheme~1 (component
states at inspection times), the likelihood for observation~$i$ takes
the form
\[
  L_i(\blam) = f_T(t_i;\blam) \prod_{j=1}^{m}
    \bigl[F_j(a_{ij}^+;\lambda_j) - F_j(a_{ij}^-;\lambda_j)\bigr],
\]
where $[a_{ij}^-,\, a_{ij}^+)$ is the inspection interval containing
$T_j$ (with the constraint $a_{ij}^+ \leq t_i$ for all~$j$ since
every component has failed by system failure time).  The interval
widths add information that resolves the left-censoring ambiguity of
Scheme~0, particularly for fast components.  Writing down and
evaluating this likelihood for exponential and Weibull components is a
natural next step; the Fisher information comparison across schemes
would quantify the value of periodic inspection.

\paragraph{Non-exponential components.}
The EM framework of \Cref{prop:em-algorithm} extends to any family where
the truncated conditional mean $\E[T_j \mid T_j < t]$ has a tractable
form.  For the Weibull family, the E-step requires the incomplete gamma
function, which is computationally available but not elementary.

\paragraph{Dependent components.}
Copula-based models could accommodate shared environmental stress or
manufacturing-induced dependence among components.  The censoring
equivalence (\Cref{prop:censoring-equiv}) continues to hold under
dependence---the system observation still comprises one exact and
$m-1$ left-censored component lifetimes---but the conditional
distributions in the E-step would no longer factor across components.


\bibliographystyle{plainnat}
\bibliography{../references}

% Appendix: Deferred Proofs
\appendix
\section{Deferred Proofs}
\label{sec:appendix-proofs}

\begin{proof}[Proof of \Cref{lem:ie}]
By induction on~$m$.

\emph{Base case} ($m = 0$).  The empty product is~$1$.  The only subset of
the empty set is $S = \varnothing$, contributing $(-1)^0 e^0 = 1$.

\emph{Inductive step.}  Assume the identity holds for $m-1$ factors.
Then
\begin{align*}
  \prod_{j=1}^m \bigl(1 - e^{-\lambda_j t}\bigr)
  &= \bigl(1 - e^{-\lambda_m t}\bigr)
     \prod_{j=1}^{m-1} \bigl(1 - e^{-\lambda_j t}\bigr) \\
  &= \bigl(1 - e^{-\lambda_m t}\bigr)
     \sum_{S \subseteq \{1,\ldots,m-1\}} (-1)^{|S|}\, e^{-\lambda(S)\, t}
     \qquad\text{[by the induction hypothesis]} \\
  &= \sum_{S \subseteq \{1,\ldots,m-1\}} (-1)^{|S|}\, e^{-\lambda(S)\, t}
     \;-\; \sum_{S \subseteq \{1,\ldots,m-1\}} (-1)^{|S|}\,
            e^{-(\lambda(S) + \lambda_m)\, t}.
\end{align*}
The first sum ranges over subsets of $\{1,\ldots,m\}$ not containing~$m$.
In the second sum, writing $S' = S \cup \{m\}$ gives
$(-1)^{|S|} = (-1)^{|S'|-1} = -(-1)^{|S'|}$, so the leading negative
sign cancels:
\[
  = \sum_{S' \subseteq \{1,\ldots,m\}} (-1)^{|S'|}\, e^{-\lambda(S')\, t}.
  \qedhere
\]
\end{proof}

\begin{proof}[Proof of \Cref{lem:density-decomp}]
Since $T = \max(T_1, \ldots, T_m)$ and the $T_j$ are independent and
continuous, ties have probability zero, and exactly one component achieves
the maximum.  The joint event $\{T \in dt,\; J = j\}$ requires
$T_j \in dt$ and $T_i < t$ for all $i \neq j$.  By independence,
\[
  P(T \in dt,\; J = j) = f_j(t) \prod_{i \neq j} F_i(t)\, dt
  = w_j(t; \blam)\, dt.
\]
Summing over $j = 1, \ldots, m$ gives the marginal density
$f_T(t; \blam) = \sum_j w_j(t; \blam)$, and dividing yields the
conditional probability~\eqref{eq:posterior-J}.
\end{proof}


\end{document}
